\clearpage
\begin{examplebox}
\textbf{\textcolor{CExample}{例 1｜典型高考题：两段直线，求参数范围}}
\begin{center}\begin{tikzpicture}[font=\small,>=stealth,line join=round]
\draw[gray!40,rounded corners=3pt] (0,0) rectangle (14.2,12.3);
\node at (7.1,11.7) {例1：两端边界为何一个不取、一个可取};
\begin{scope}[shift={(3.0,7.8)},x=0.48cm,y=0.48cm]\draw[->,gray] (-2,0)--(2.3,0) node[right]{$x$};\draw[->,gray] (0,-1)--(0,7.3) node[left]{$y$};\begin{scope}\clip (-2,-1) rectangle (2,7);\draw[blue,thick,domain=-2:1,samples=80] plot(\x,{0*\x+2});\draw[orange!85!black,thick,domain=1.0001:2,samples=100] plot(\x,{3*\x+1});\end{scope}\draw[draw=blue,fill=blue] (1,2) circle(2pt);\draw[draw=orange!85!black,fill=white] (1,4) circle(2pt);\end{scope}\node at (3.0, 6.75) {$a=1$：左段平：失败};\begin{scope}[shift={(10.0,7.8)},x=0.48cm,y=0.48cm]\draw[->,gray] (-2,0)--(2.3,0) node[right]{$x$};\draw[->,gray] (0,-1)--(0,7.3) node[left]{$y$};\begin{scope}\clip (-2,-1) rectangle (2,7);\draw[blue,thick,domain=-2:1,samples=80] plot(\x,{1*\x+2});\draw[orange!85!black,thick,domain=1.0001:2,samples=100] plot(\x,{3*\x+1});\end{scope}\draw[draw=blue,fill=blue] (1,3) circle(2pt);\draw[draw=orange!85!black,fill=white] (1,4) circle(2pt);\end{scope}\node at (10.0, 6.75) {$a=2$：向上跳：通过};\begin{scope}[shift={(3.0,2.5)},x=0.48cm,y=0.48cm]\draw[->,gray] (-2,0)--(2.3,0) node[right]{$x$};\draw[->,gray] (0,-1)--(0,7.3) node[left]{$y$};\begin{scope}\clip (-2,-1) rectangle (2,7);\draw[blue,thick,domain=-2:1,samples=80] plot(\x,{2*\x+2});\draw[orange!85!black,thick,domain=1.0001:2,samples=100] plot(\x,{3*\x+1});\end{scope}\draw[draw=blue,fill=blue] (1,4) circle(2pt);\end{scope}\node at (3.0, 1.45) {$a=3$：接平：通过};\begin{scope}[shift={(10.0,2.5)},x=0.48cm,y=0.48cm]\draw[->,gray] (-2,0)--(2.3,0) node[right]{$x$};\draw[->,gray] (0,-1)--(0,7.3) node[left]{$y$};\begin{scope}\clip (-2,-1) rectangle (2,7);\draw[blue,thick,domain=-2:1,samples=80] plot(\x,{3*\x+2});\draw[orange!85!black,thick,domain=1.0001:2,samples=100] plot(\x,{3*\x+1});\end{scope}\draw[draw=blue,fill=blue] (1,5) circle(2pt);\draw[draw=orange!85!black,fill=white] (1,4) circle(2pt);\end{scope}\node at (10.0, 1.45) {$a=4$：向下跳：失败};
\node[text width=12.8cm,align=center] at (7.1,.65) {左斜率 $a-1>0$，接口 $a+1\le4$；共同给出 $1<a\le3$。各小图横纵等比例。};
\end{tikzpicture}\end{center}



\[
 f(x)=\begin{cases}(a-1)x+2,&x\le1,\\3x+1,&x>1.\end{cases}
\]
问：$f(x)$ 在 $\mathbb R$ 上\textbf{严格递增}，$a$ 取什么值？

三关完整展开：
\[
\begin{aligned}
\text{左段：}\quad &a-1>0\quad\Longrightarrow\quad a>1;\\
\text{右段：}\quad &3>0\quad\text{（自动通过）};\\
\text{接口：}\quad &g(1)=(a-1)\cdot1+2=a+1,\quad L=3\cdot1+1=4,\\
&a+1\le4\quad\Longrightarrow\quad a\le3.
\end{aligned}
\]
两关合起来：\[\boxed{1<a\le3}.\]
\textbf{检查端点：}$a=1$ 时左段变成水平线，不再严格递增，必须去掉；$a=3$ 时两边的高度正好相接，\emph{仍严格递增}，所以可以保留。

\medskip
\end{examplebox}
\clearpage
\begin{examplebox}
\textbf{\textcolor{CExample}{例 2｜整体严格递减：不等号要反向}}
\begin{center}\begin{tikzpicture}[font=\small,>=stealth,line join=round]
\draw[gray!40,rounded corners=3pt] (0,0) rectangle (14.2,8.0);
\node at (7.1,7.4) {例2：递减必须左高右低，接口取等也可};
\begin{scope}[shift={(2.0,2.7)},x=0.55cm,y=0.55cm]\draw[->,gray] (-2,0)--(2.3,0) node[right]{$x$};\draw[->,gray] (0,0)--(0,7.3) node[above]{$y$};\begin{scope}\clip (-2,0) rectangle (2,7);\draw[blue,thick,domain=-2:0,samples=80] plot(\x,{-2*\x+2});\draw[orange!85!black,thick,domain=0.0001:2,samples=100] plot(\x,{-\x+3});\end{scope}\draw[draw=blue,fill=blue] (0,2) circle(2pt);\draw[draw=orange!85!black,fill=white] (0,3) circle(2pt);\end{scope}\node at (2.0,1.3) {$a=2$：失败};\begin{scope}[shift={(6.6,2.7)},x=0.55cm,y=0.55cm]\draw[->,gray] (-2,0)--(2.3,0) node[right]{$x$};\draw[->,gray] (0,0)--(0,7.3) node[above]{$y$};\begin{scope}\clip (-2,0) rectangle (2,7);\draw[blue,thick,domain=-2:0,samples=80] plot(\x,{-2*\x+3});\draw[orange!85!black,thick,domain=0.0001:2,samples=100] plot(\x,{-\x+3});\end{scope}\draw[draw=blue,fill=blue] (0,3) circle(2pt);\end{scope}\node at (6.6,1.3) {$a=3$：通过};\begin{scope}[shift={(11.2,2.7)},x=0.55cm,y=0.55cm]\draw[->,gray] (-2,0)--(2.3,0) node[right]{$x$};\draw[->,gray] (0,0)--(0,7.3) node[above]{$y$};\begin{scope}\clip (-2,0) rectangle (2,7);\draw[blue,thick,domain=-2:0,samples=80] plot(\x,{-2*\x+4});\draw[orange!85!black,thick,domain=0.0001:2,samples=100] plot(\x,{-\x+3});\end{scope}\draw[draw=blue,fill=blue] (0,4) circle(2pt);\draw[draw=orange!85!black,fill=white] (0,3) circle(2pt);\end{scope}\node at (11.2,1.3) {$a=4$：通过};
\node[text width=12.8cm,align=center] at (7.1,.65) {蓝段 $-2x+a$；橙段 $-x+3$。两段都减，整体严格递减恰好要求 $a\ge3$。};
\end{tikzpicture}\end{center}


\[
 f(x)=\begin{cases}-2x+a,&x\le0,\\-x+3,&x>0.\end{cases}
\]
两段斜率 $-2,-1$ 都是负数，所以两段各自严格递减。接口处左边的函数值是 $a$，右边刚开始接近 $3$。递减需要左边不低于右边：
\[
 a\ge3.\qquad\boxed{a\in[3,+\infty)}.
\]
\textbf{检查：}$a=3$ 也可以，不能误写成 $a>3$。

\medskip
\end{examplebox}
\clearpage
\begin{examplebox}
\textbf{\textcolor{CExample}{例 3｜“跳起来”也能严格递增}}
\begin{center}\begin{tikzpicture}[font=\small,>=stealth,line join=round]
\draw[gray!40,rounded corners=3pt] (0,0) rectangle (14.2,7.4);
\node at (7.1,6.800000000000001) {两段斜率都为正，交接方向决定整体是否递增};
\begin{scope}[shift={(2.1,3.1)},x=0.65cm,y=0.65cm]\draw[->,gray] (-2,0)--(2.3,0) node[right]{$x$};\draw[->,gray] (0,-2)--(0,4.3) node[above]{$y$};\begin{scope}\clip (-2,-2) rectangle (2,4);\draw[blue,thick,domain=-2:0,samples=80] plot(\x,{\x});\draw[orange!85!black,thick,domain=0.0001:2,samples=100] plot(\x,{\x+2});\end{scope}\draw[draw=blue,fill=blue] (0,0) circle(2pt);\draw[draw=orange!85!black,fill=white] (0,2) circle(2pt);\end{scope}\node at (2.1,1.35) {向上跳：严格增};\begin{scope}[shift={(6.699999999999999,3.1)},x=0.65cm,y=0.65cm]\draw[->,gray] (-2,0)--(2.3,0) node[right]{$x$};\draw[->,gray] (0,-2)--(0,4.3) node[above]{$y$};\begin{scope}\clip (-2,-2) rectangle (2,4);\draw[blue,thick,domain=-2:0,samples=80] plot(\x,{\x});\draw[orange!85!black,thick,domain=0.0001:2,samples=100] plot(\x,{\x});\end{scope}\draw[draw=blue,fill=blue] (0,0) circle(2pt);\draw[draw=blue,fill=blue] (0,0) circle(2pt);\end{scope}\node at (6.699999999999999,1.35) {接平：仍严格增};\begin{scope}[shift={(11.299999999999999,3.1)},x=0.65cm,y=0.65cm]\draw[->,gray] (-2,0)--(2.3,0) node[right]{$x$};\draw[->,gray] (0,-2)--(0,4.3) node[above]{$y$};\begin{scope}\clip (-2,-2) rectangle (2,4);\draw[blue,thick,domain=-2:0,samples=80] plot(\x,{\x+2});\draw[orange!85!black,thick,domain=0.0001:2,samples=100] plot(\x,{\x});\end{scope}\draw[draw=blue,fill=blue] (0,2) circle(2pt);\draw[draw=orange!85!black,fill=white] (0,0) circle(2pt);\end{scope}\node at (11.299999999999999,1.35) {向下跳：不递增};
\node[text width=12.8cm,align=center] at (7.1,.65) {蓝实心取，橙空心不取；接平时端点重合，按实心取。};
\end{tikzpicture}\end{center}


\[
 f(x)=\begin{cases}x,&x\le0,\\x+2,&x>0.\end{cases}
\]
左右斜率都是 $1>0$；左值 $g(0)=0$，右边刚开始的高度是 $2$，满足 $0<2$。因此整个函数仍然\textbf{严格递增}，虽然在 $0$ 处图象断开。

相反，若两段换成 $x+2\ (x\le0)$ 和 $x\ (x>0)$，则左值 $2$ 大于右边刚开始的 $0$，不满足接口条件，\textbf{整体不递增}。

\medskip
\end{examplebox}
\clearpage
\begin{examplebox}
\textbf{\textcolor{CExample}{例 4｜指数与对数拼接：三关一次写完}}
\begin{center}\begin{tikzpicture}[font=\small,>=stealth,line join=round]
\draw[gray!40,rounded corners=3pt] (0,0) rectangle (14.2,7.6);
\node at (7.1,7.0) {例4：指数与对数都增，还要比较两端高度};
\begin{scope}[shift={(2.0,2.6)},x=0.68cm,y=0.68cm]\draw[->,gray] (-2,0)--(2.3,0) node[right]{$x$};\draw[->,gray] (0,0)--(0,4.8) node[above]{$y$};\begin{scope}\clip (-2,0) rectangle (2,4.5);\draw[blue,thick,domain=-2:0,samples=80] plot(\x,{2*pow(3,\x)});\draw[orange!85!black,thick,domain=0.0001:2,samples=100] plot(\x,{ln(1+\x)+1});\end{scope}\draw[draw=blue,fill=blue] (0,2) circle(2pt);\draw[draw=orange!85!black,fill=white] (0,1) circle(2pt);\end{scope}\node at (2.0,1.25) {$a=2,b=1$：失败};\begin{scope}[shift={(6.6,2.6)},x=0.68cm,y=0.68cm]\draw[->,gray] (-2,0)--(2.3,0) node[right]{$x$};\draw[->,gray] (0,0)--(0,4.8) node[above]{$y$};\begin{scope}\clip (-2,0) rectangle (2,4.5);\draw[blue,thick,domain=-2:0,samples=80] plot(\x,{2*pow(3,\x)});\draw[orange!85!black,thick,domain=0.0001:2,samples=100] plot(\x,{ln(1+\x)+2});\end{scope}\draw[draw=blue,fill=blue] (0,2) circle(2pt);\end{scope}\node at (6.6,1.25) {$a=2,b=2$：通过};\begin{scope}[shift={(11.2,2.6)},x=0.68cm,y=0.68cm]\draw[->,gray] (-2,0)--(2.3,0) node[right]{$x$};\draw[->,gray] (0,0)--(0,4.8) node[above]{$y$};\begin{scope}\clip (-2,0) rectangle (2,4.5);\draw[blue,thick,domain=-2:0,samples=80] plot(\x,{2*pow(3,\x)});\draw[orange!85!black,thick,domain=0.0001:2,samples=100] plot(\x,{ln(1+\x)+3});\end{scope}\draw[draw=blue,fill=blue] (0,2) circle(2pt);\draw[draw=orange!85!black,fill=white] (0,3) circle(2pt);\end{scope}\node at (11.2,1.25) {$a=2,b=3$：通过};
\node[text width=12.8cm,align=center] at (7.1,.65) {蓝段 $a3^x$，橙段 $\ln(1+x)+b$；须 $a>0$ 且 $a\le b$。此图固定 $a=2$。};
\end{tikzpicture}\end{center}


\[
 f(x)=\begin{cases}a\,3^x,&x\le0,\\\ln(1+x)+b,&x>0.\end{cases}
\]
题目要求整体严格递增。

\textbf{左段：}$3^x$ 本来严格递增，乘上 $a$ 后仍严格递增须 $a>0$。

\textbf{右段：}$1+x$ 随 $x$ 增大而增大，$\ln(1+x)$ 也随之严格增大；加上 $b$ 只会整体上下移动，不改变增减方向。

\textbf{接口：}左值 $g(0)=a$，右边靠近 $0$ 的高度是 $\ln 1+b=b$，故 $a\le b$。

答案：\[\boxed{a>0,\qquad b\ge a}.\]

\medskip
\end{examplebox}
\clearpage
\begin{examplebox}
\textbf{\textcolor{CExample}{例 5｜等号写在哪一侧，先看清楚}}
\begin{center}\begin{tikzpicture}[font=\small,>=stealth,line join=round]
\draw[gray!40,rounded corners=3pt] (0,0) rectangle (14.2,10.3);
\node at (7.1,9.700000000000001) {例5：等号在右侧，实心点也属于右侧};
\begin{scope}[shift={(3.1,4.9)},x=0.9cm,y=0.9cm]\draw[->,gray] (-2,0)--(2.3,0) node[right]{$x$};\draw[->,gray] (0,-4)--(0,4.3) node[above]{$y$};\begin{scope}\clip (-2,-4) rectangle (2,4);\draw[blue,thick,domain=-2:0,samples=80] plot(\x,{\x-2});\draw[orange!85!black,thick,domain=0.0001:2,samples=100] plot(\x,{\x+1});\end{scope}\draw[draw=blue,fill=white] (0,-2) circle(2pt);\draw[draw=orange!85!black,fill=orange!85!black] (0,1) circle(2pt);\end{scope}\node[align=left,anchor=west] at (7.5,5.1) {$x<0$：蓝段 $x-2$\\[12pt]$x\ge0$：橙段 $x+1$\\[12pt]空心：左极限 $-2$ 不取\\[9pt]实心：真正的 $f(0)=1$\\[12pt]$-2\le1$，整体严格增};
\node[text width=12.8cm,align=center] at (7.1,.65) {先看端点归属，再比较左右高度；不能把未取的蓝点当作 $f(0)$。};
\end{tikzpicture}\end{center}


\[
 f(x)=\begin{cases}x-2,&x<0,\\x+1,&x\ge0.\end{cases}
\]
这里 $f(0)=1$，\emph{等号属于右段}。左边靠近 $0$ 时函数值接近 $-2$；右边从 $1$ 开始，函数仍严格递增。

判法：左右各自递增，且\textbf{左边最后接近的高度} $-2$ \textbf{不高于} $f(0)=1$。不能照搬“比较 $g(0)$ 与右边高度”的形式，因为本题 $x=0$ 不属于左段。

\medskip
\end{examplebox}
\clearpage
\begin{examplebox}
\textbf{\textcolor{CExample}{例 6｜公式不能直接用时，回到最朴素的比较}}
\begin{center}\begin{tikzpicture}[font=\small,>=stealth,line join=round]
\draw[gray!40,rounded corners=3pt] (0,0) rectangle (14.2,9.0);
\node at (7.1,8.4) {例6：右段无有限起始高度，直接找逆序点};
\begin{scope}[shift={(3.4,6.8)},x=0.55cm,y=0.55cm]\draw[->,gray] (-3,0)--(3.3,0) node[right]{$x$};\draw[->,gray] (0,-10.5)--(0,1.8) node[above]{$y$};\begin{scope}\clip (-3,-10.5) rectangle (3,1.5);\draw[blue,thick,domain=-3:0,samples=80] plot(\x,{\x});\draw[orange!85!black,thick,domain=0.095:3,samples=100] plot(\x,{-1/\x});\end{scope}\draw[draw=blue,fill=blue] (0,0) circle(2pt);\end{scope}\begin{scope}[shift={(3.4,6.8)},x=.55cm,y=.55cm]\fill[red!70!black] (.1,-10) circle(2.5pt);\draw[red!70!black,->] (1.3,-8.7)--(.1,-10);\end{scope}\node[align=left,anchor=west] at (7.2,4.5) {$f(0)=0$，$f(0.1)=-10$\\[12pt]$0<0.1$，函数值却变小\\[12pt]右段在 $x>0$ 上仍严格增\\[10pt]$x\to0^+$ 时 $-1/x\to-\infty$\\[12pt]没有有限的接口高度 $L$};
\node[text width=12.8cm,align=center] at (7.1,.65) {蓝段 $x$，橙段 $-1/x$；橙色曲线在 $x=0$ 无端点，图中下方只是视窗裁切。};
\end{tikzpicture}\end{center}


\[
 f(x)=\begin{cases}x,&x\le0,\\-\dfrac1x,&x>0.\end{cases}
\]
左段严格递增；右段 $-1/x$ 在 $x>0$ 也严格递增。但在 $x=0$，函数值是 $0$，到 $x=0.1$ 时却变成 $-10$，因此整体不递增。

这里右段靠近 $0$ 时函数值不停往更小的方向走，\textbf{没有有限的起始高度 $L$}。所以不能套用 $g(0)\le L$ 的有限数公式；直接比较 $f(0)$ 与 $f(0.1)$，马上就知道答案。
\end{examplebox}
